\exercisegroup

\begin{enumerate}
\def\labelenumi{\arabic{enumi})}
\item
  Let A be a convex set. Then, a point \(x \in A\) is extremal in A if, and only if, for any subset B of A, the statement x belongs to the convex envelope of B, implies that \(x \in B\) .
\item
  With the notation of II, p.~74, exerc. 3, let G be the vector subspace of E generated by K ∪ \{λ\}. Show that, in G, the point λ is an extremal point of the closed convex envelope of K, but that λ does not belong to K (cf.~II, p.~25, corollary).
\end{enumerate}

¶ 3) Let A be a convex set in a vector space E, and let x be a point of A. We call the set formed by x, and the \(y \neq x\) in A such that the line passing through x and y contains an open segment, that is contained in A and contains x, the facet of x in A. The internal points relative to the linear variety generated by A (II, p.~26) (resp. the extremal points) of A are the points whose facet in A is equal to A (resp. is a single point).

\begin{enumerate}
\def\labelenumi{\alph{enumi})}
\item
  Show that the facet \(\mathrm{F}_x\) of a point \(x \in \mathbf{A}\) is the largest convex set \(\mathbf{B} \subset \mathbf{A}\) such that \(x\) is an internal point of \(\mathbf{B}\) (relative to the linear variety generated by \(\mathbf{B}\) ).
\item
  For every point \(y \in \mathbf{F}_x\) , the facet \(\mathbf{F}_y\) of \(y\) in \(\mathbf{A}\) is identical with the facet of \(y\) in \(\mathbf{F}_x\) . In order that \(\mathbf{F}_y = \mathbf{F}_x\) , it is necessary and sufficient that \(y\) is an internal point of \(\mathbf{F}_x\) (relative to the linear variety generated by \(\mathbf{F}_x\) ). Deduce that, if \(\mathbf{F}_x\) is finite dimensional, and if \(y\) is a non-internal point of \(\mathbf{F}_x\) (relative to the linear variety generated by \(\mathbf{F}_x\) ), then the dimension of \(\mathbf{F}_y\) is strictly less than that of \(\mathbf{F}_x\) .
\item
  A linear variety V in E which meets A and is such that for every \(x \in A \cap V\) , every open segment contained in A and containing x, is necessarily contained in V, is called a support variety of A. Show that, for all \(x \in A\) , the linear variety M generated by the facet \(F_{x}\) of x in A is the smallest support variety of A which contains x, and that \(M \cap A = F_{x}\) . For every support variety V of A, the intersection \(V \cap A\) is the facet in A of each of its internal points (relative to the linear variety generated by \(V \cap A\) ).
\end{enumerate}

\(d\) ) Let \(\mathbf{A}\) and \(\mathbf{B}\) be two convex sets in \(\mathbf{E}\) . For every point \(x\in \mathbf{A}\cap \mathbf{B}\) , the facet of \(x\) in \(\mathbf{A}\cap \mathbf{B}\) is the intersection of the facets of \(x\) in \(\mathbf{A}\) and in \(\mathbf{B}\) .

\begin{enumerate}
\def\labelenumi{\alph{enumi})}
\setcounter{enumi}{4}
\item
  Let B be a closed convex set in a Hausdorff topological vector space E, and let B contain a closed linear variety M of finite codimension n; then every facet in B of a point of B contains a closed linear variety of codimension n (II, p.~67, exerc. 14, d)). If A is a convex set then the facet in A ∩ B of a point x in A ∩ B is of finite dimension if, and only if, the facet of x in A is of finite dimension: further if p and q are the dimension of the facet of x in A and of the facet of x in A ∩ B, then \(p \leqslant q + n\) . In particular if \(x \in A \cap B\) is an extremal point of A ∩ B, then its facet in A is of dimension \(\leqslant n\) .
\item
  Deduce from e) that if A is compact, and V is a closed linear variety in E, of finite codimension n, then every extremal point of \(V \cap A\) is a linear combination of at most \(n + 1\) extremal points of A.
\end{enumerate}

\begin{enumerate}
\def\labelenumi{\arabic{enumi})}
\setcounter{enumi}{3}
\item
  In the plane \(\mathbf{R}^2\) , consider the convex set \(\mathbf{A}\) formed by the points \((\xi, \eta)\) satisfying \(-1 \leqslant \xi \leqslant 1\) , \(-1 - \sqrt{1 - \xi^2} \leqslant \eta \leqslant 1 + \sqrt{1 - \xi^2}\) . Show that there exist frontier points of \(\mathbf{A}\) for which the facet is distinct from the intersection of \(\mathbf{A}\) and of the lines of support of \(\mathbf{A}\) passing through this point.
\item
  In the Banach space \(l^{\infty}(\mathbf{N})\) of bounded sequences \(x = (\xi_n)\) of real numbers, let \(\mathbf{A}\) be the closed convex set defined by the inequalities \(-1 / n \leqslant \xi_n \leqslant 1\) for \(n \geqslant 1\) and \(-1 \leqslant \xi_0 \leqslant 1\) . Show that \(\mathbf{A}\) has a non-empty interior, that the origin is a frontier point of \(\mathbf{A}\) and that the facet of 0 in \(\mathbf{A}\) is not closed. If we give to \(\mathbf{A}\) the topology induced by that of the product space \(\mathbf{R}^{\mathbf{N}}\) , show that \(\mathbf{A}\) is compact but that the facet of 0 in \(\mathbf{A}\) is not closed in \(\mathbf{A}\) .
\item
  Let E, \(E'\) be two vector spaces in separating duality, and A be a convex set in E containing 0 and closed for \(\sigma(\mathrm{E}, \mathrm{E}')\) . For all \(a \in A\) , the set \(F_{a}'\) of points \(x' \in A^{\circ}\) such that \(\langle a, x' \rangle = -1\) is a closed (for \(\sigma(\mathrm{E}', \mathrm{E})\) ), convex set of \(A^{\circ}\) . Show that \(F_{a}'\) is the facet in \(A^{\circ}\) of each of the internal points of \(F_{a}'\) relative to the linear variety generated by \(F_{a}'\) . We say that \(F_{a}'\) is the dual facet of a in \(A^{\circ}\) . If \(F_{a}\) is the facet of a in A, show that \(F_{a}'\) is also the dual facet in \(A^{\circ}\) of each of the internal points of \(F_{a}\) relative to the linear variety generated by \(F_{a}\) ; further, if A is identified with \(A^{\circ\circ}\) , the dual facet in A of each internal point of \(F_{a}^{\prime}\) relative to the linear variety generated by \(F_{a}^{\prime}\) , contains \(F_{a}\) . When \(F_{a}^{\prime}\) is not empty (which is always the case when E is finite-dimensional and \(a \neq 0\) , cf.~II, p.~78, exerc. 13), we say that each of \(F_{a}\) , \(F_{a}^{\prime}\) is the dual facet of the other.
\end{enumerate}

We say that a point \(a \in \mathbf{A}\) is smooth point of \(\mathbf{A}\) if \(\mathbf{F}_a'\) is a single point (in other words if there exists one closed hyperplane of support of \(\mathbf{A}\) passing through \(a\) ); we say that \(a\) is a point of strict convexity (or is an exposed point) if there exists a closed hyperplane \(\mathbf{H}\) supporting \(\mathbf{A}\) so that \(\mathbf{H} \cap \mathbf{A} = \{a\}\) ; this is the same as saying that there exists an internal point of \(\mathbf{F}_a'\) (relative to a linear variety generated by \(\mathbf{F}_a'\) ) which is a smooth point of \(\mathbf{A}^\circ\) .

\begin{enumerate}
\def\labelenumi{\arabic{enumi})}
\setcounter{enumi}{6}
\tightlist
\item
  Let \(\mathbf{E}\) be a vector space of finite dimension \(n\) and let \(\mathbf{A}\) be a closed convex set in \(\mathbf{E}\) of which 0 is an interior point.
\end{enumerate}

\begin{enumerate}
\def\labelenumi{\alph{enumi})}
\item
  Let F and \(F'\) be two dual facets of A and \(A^{\circ}\) (exerc. 6); if F is of dimension p and \(F'\) of dimension q, then show that \(p + q \leqslant n - 1\) . For every frontier point x of A, the dimension of the facet of x in A is called the order of x, and the dimension of its dual facet in \(A^{\circ}\) is called the class of x. The order (resp. the class) of a facet F of A is by definition the order (resp. the class) of one of the internal points of F relative to the linear variety generated by F. An extremal point of A is a point of order 0; a smooth point of A (exerc. 6) is a point of class 0.
\item
  A frontier point of A of class n - 1 (and hence of order 0) is called a vertex of A. Show that the set of vertices of A is enumerable (consider the set of dual facets of the vertices of A, GT, VI, § 2, exerc. 12).
\item
  Let F be a p-dimensional facet of A, and M a linear variety of dimension n - p, which meets F in the single point a, such that a is an internal point of F and which contains an interior point of A. Show that, if V is a support hyperplane of \(M \cap A\) in M, that passes through a, then the hyperplane H generated (in E) by \(F \cup V\) is a support hyperplane of A.
\item
  We say that a facet F of A of order p and of class q is an ultrafacet if \(p + q = n - 1\) ; the dual facet is then also an ultrafacet of \(A^{\circ}\) . If a linear variety M of dimension n - p meets an ultrafacet F in a single point that is an internal point of F (relative to the linear variety generated by F), show that this point is a vertex of the convex set \(M \cap A\) , and conversely (use c). Deduce that the set of ultrafacets of order p of A is enumerable. (Identify E with \(R^{n}\) , consider the projection of A on each of the coordinate varieties of \(R^{n}\) of dimension p; if the set of ultrafacets of order p of which the projection on V is p-dimensional, is not enumerable, show that there exists a point of V which is an interior point to a non-enumerable infinity of these projections considering the points of V with rational coordinates; then use b.) Give an example of a convex set with a non-enumerable infinity of facets, each of which is not a single point nor an ultrafacet.
\item
  If all the frontier points of A are smooth, show that the mapping, which puts each point x of the frontier G of A in correspondence with the unique point of the dual facet of x, is a continuous mapping of G on the frontier of A° (cf.~TG, I, § 9.1, corollary). In what case is this mapping bijective?
\end{enumerate}

\begin{enumerate}
\def\labelenumi{\arabic{enumi})}
\setcounter{enumi}{7}
\tightlist
\item
  Let \(\mathbf{E}\) be a vector space of finite dimension \(n\) and \(\mathbf{A}\) be a compact convex set in \(\mathbf{E}\) .
\end{enumerate}

\begin{enumerate}
\def\labelenumi{\alph{enumi})}
\item
  Let H be a hyperplane in E. Show that in an open half-space determined by H and containing at least one point of A, there exists a point of strict convexity of A (II, p.~87, exerc. 6). (Consider, in H, a closed euclidean ball C of dimension n - 1 and of sufficiently large radius that contains H ∩ A, then the euclidean balls B of dimension n and of larger radius containing A and such that B ∩ H = C.)
\item
  Show that A is the closed convex envelope of the set of points of strict convexity (use a)). c) Show that every extremal point of A is a cluster point of the set of points of A of strict convexity. (Using b) and the exerc. 9, a) of II, p.~66, note that an extremal point is the limit of a sequence of points of the form \(\sum_{i=0}^{n} \lambda_{im} x_{im}\) , where \(\lambda_{im} \geqslant 0\) , \(\sum_{i=0}^{n} \lambda_{im} = 1\) and the \(x_{im}\) are points of strict convexity of A; next use the compactness of A.)
\end{enumerate}

\begin{enumerate}
\def\labelenumi{\arabic{enumi})}
\setcounter{enumi}{8}
\item
  Show that in the product space \(\mathbf{E} = \mathbf{R}^{\mathbf{N}}\) , the cube \(\mathbf{I}^{\mathbf{N}}\) , where \(\mathbf{I} = [0,1]\) , is a compact convex set with no point of strict convexity.
\item
  In the space \(\mathbf{R}^2\) , show that the set of extremal points of a closed convex set \(\mathbf{A}\) is closed (show that the set of points of A whose facet in A is of dimension 1 form an open set relative to the frontier of A).
\item
  \begin{enumerate}
  \def\labelenumii{\alph{enumii})}
  \tightlist
  \item
    In the space \(\mathbf{R}^3\) , consider the compact convex set A which is the convex envelope of the union of the circle \(\zeta = 0\) , \(\xi^2 + \eta^2 - 2\xi = 0\) and the two points \((0, 0, 1)\) and \((0, 0, -1)\) . Show that the set of extremal points of A is not closed in A.
  \end{enumerate}
\end{enumerate}

\begin{enumerate}
\def\labelenumi{\alph{enumi})}
\setcounter{enumi}{1}
\tightlist
\item
  Let A be a metrisable compact convex set in a Hausdorff topological vector space E. Show that the set of extremal points of A is the intersection of a sequence of open sets in A. (If d is a distance defining the topology of A, then for each integer n consider the set of points \(x = \frac{1}{2}(y + z)\) , where y, z are in A and \(d(y, z) \geqslant 1/n\) .)
\end{enumerate}

\begin{enumerate}
\def\labelenumi{\arabic{enumi})}
\setcounter{enumi}{11}
\tightlist
\item
  In the Banach space \(l^{\infty}(\mathbf{N})\) , let \(e_n\) be the sequence all of whose terms are zero except the \(n\) -th which is 1. Let \(\mathbf{A}\) be the convex closed envelope of the set formed from 0 and the points \(\mathbf{e}_n / (n + 1)\) ( \(n \geqslant 0\) ). Show that \(\mathbf{A}\) is compact but that it is not identical with the convex envelope of the set of its extremal points.
\end{enumerate}

* 13) In the Hilbert space \(l^2(\mathbf{N})\) , let \(\mathbf{A}\) be the set of points \(x = (\xi_n)\) such that we have \(\sum_{n} 2^{2n}\xi_n^2 \leqslant 1\) . Show that \(\mathbf{A}\) is convex, compact and that it is the closure of the set of its extremal points.

\begin{enumerate}
\def\labelenumi{\arabic{enumi})}
\setcounter{enumi}{13}
\tightlist
\item
  Let \(\mathbf{E}\) be a closed vector subspace of the Banach space \(l^{\infty}(\mathbf{N})\) , formed of the sequences \(x = (\xi_n)\) such that \(\lim_{n\to \infty}\xi_n = 0\) .
\end{enumerate}

\begin{enumerate}
\def\labelenumi{\alph{enumi})}
\tightlist
\item
  Show that, in the Banach space E, the closed unit ball B does not have any extremal points. b) Let \(u\) be the continuous linear form \((\xi_n) \mapsto \sum_n 2^{-n}\xi_n\) on E. Show that there does not exist any support hyperplane of B that is parallel to the closed hyperplane with the equation \(u(x)=0\) .
\end{enumerate}

\begin{enumerate}
\def\labelenumi{\arabic{enumi})}
\setcounter{enumi}{14}
\tightlist
\item
  Let \(\mathbf{A}\) be a compact set in the normed space \(\mathbf{E}\) .
\end{enumerate}

\begin{enumerate}
\def\labelenumi{\alph{enumi})}
\item
  Show that the distance apart of two parallel support hyperplanes of A is at most equal to the diameter \(\delta\) of A.
\item
  Show that there exist pairs of points \((a, b)\) of A such that \(\|a - b\| = \delta\) ; for such a pair of points, there exist two parallel support hyperplanes of A passing respectively through a and b and whose distance apart is \(\delta\) (consider the closed ball of centre a and radius \(\delta\) ).
\end{enumerate}

\begin{enumerate}
\def\labelenumi{\arabic{enumi})}
\setcounter{enumi}{15}
\tightlist
\item
  \begin{enumerate}
  \def\labelenumii{\alph{enumii})}
  \tightlist
  \item
    Let A be an n-dimensional compact convex set in the space \(R^{n}\) , normed with the Euclidean norm; for every \(z \in S_{n-1}\) , denote by \(\rho(z)\) the upper bound of lengths of segments parallel to the vector z and contained in A. Show that there exist two points u, v of A such that the segment with end points u, v is parallel to z and is of length \(\rho(z)\) ; deduce that there exist two support hyperplanes of A that are parallel and pass respectively through u and v (consider the set \(A' = A + \rho(z)z\) , and separate the sets A and \(A'\) by a hyperplane).
  \end{enumerate}
\end{enumerate}

\begin{enumerate}
\def\labelenumi{\alph{enumi})}
\setcounter{enumi}{1}
\tightlist
\item
  Let d be the lower bound of the distances between two parallel support hyperplanes of A; show that there exist two points a, b of A such that \(\|a - b\| = d\) , and that the hyperplanes passing respectively through a and b and orthogonal to a - b, are support hyperplanes of A (use a)).
\end{enumerate}

\begin{enumerate}
\def\labelenumi{\arabic{enumi})}
\setcounter{enumi}{16}
\item
  In the space \(l^{\infty}(\mathbf{N})\) , let \(\mathbf{A}\) be the compact convex set defined in the exerc. 12 of II, p.~89 and let \(\mathbf{E}\) be the closed vector subspace of \(l^{\infty}(\mathbf{N})\) generated by \(\mathbf{A}\) . Show that the lower bound of the distance between two parallel closed support hyperplanes of \(\mathbf{A}\) in the space \(\mathbf{E}\) is equal to 0, even though \(\mathbf{A}\) is not contained in a closed hyperplane of \(\mathbf{E}\) .
\item
  In a Hausdorff locally convex space E, let \((\mathbf{K}_{\alpha})_{\alpha\in\mathbf{I}}\) be a decreasing directed family of convex sets that are compact and non-empty. For all \(\alpha\in I\) , denote the set of extremal points of \(K_{\alpha}\) by \(A_{\alpha}\) , and by \(F_{\alpha}\) the closure of the union of the \(A_{\beta}\) for \(\beta\geqslant\alpha\) , so that \((\mathrm{F}_{\alpha})\) is a decreasing directed family of compact sets. Let A be the intersection (non-empty) of the \(F_{\alpha}\) , and K the intersection (non-empty) of the \(K_{\alpha}\) . Show that K is the closed convex envelope of A. (If f is a continuous linear form on E, and \(x_{\alpha}\) a point of \(F_{\alpha}\) where f' attains its maximum in \(F_{\alpha}\) , show that \(f(y)\leqslant f(x_{\alpha})\) for all \(y\in K\) ; then take a cluster point of the family \((x_{\alpha})\) following the filter of sections of I.)
\end{enumerate}

¶ 19) In the space \(\mathbf{R}^n\) , let \((\mathrm{K}_{\alpha})\) be a family of compact sets, in number \(\geqslant n + 1\) , and such that none of them is contained in an affine hyperplane. Suppose that for every family \((u_{\alpha})\) of affine automorphisms of \(\mathbf{R}^n\) , if any \(n + 1\) of the sets \(u_{\alpha}(\mathrm{K}_{\alpha})\) have a common point, then all the \(u_{\alpha}(\mathrm{K}_{\alpha})\) have a common point. Show that under these conditions the \(\mathbf{K}_{\alpha}\) are convex. (Suppose on the contrary that there exist \(n + 1\) points \(x_1, \ldots, x_{n+1}\) in the same set \(\mathbf{K}_{\alpha}\) and a point \(x_0\) which belongs to the convex envelope of the set of the \(x_i\) ( \(i \geqslant 1\) ) but does not belong to \(\mathbf{K}_{\alpha}\) . Note that for every index \(i \geqslant 1\) , there is an affine automorphism \(u_i\) of \(\mathbf{R}^n\) and an index \(\alpha_i\) such that \(x_0\) and the \(x_j\) of index \(j \neq i\) are extremal points of \(u_i(\mathrm{K}_{\alpha_i})\) , and show that the \(n + 2\) sets \(\mathbf{K}_{\alpha}\) and \(u_i(\mathrm{K}_{\alpha_i})\) have no points in common.)

\begin{enumerate}
\def\labelenumi{\arabic{enumi})}
\setcounter{enumi}{19}
\tightlist
\item
  Give an example of a compact convex set \(\mathbf{K}\) in \(\mathbf{R}^2\) , containing 0 and such that the cone of vertex 0 generated by \(\mathbf{K}\) is not closed in \(\mathbf{R}^2\) .
\end{enumerate}

¶ 21) a) In a Hausdorff locally convex space E, let A be a locally compact closed convex cone, that does not contain any line. Show that A is a cone of compact sole (apply prop. 2 of II, p.~55, to the vertex of A, which is an extremal point of A). Deduce that there exists a closed support hyperplane H of A, which contains the vertex s of A and is such that \(H \cap A = \{s\}\) . b) Let A, B be two closed convex cones with vertex 0 in E, that are locally compact and do not contain a line. Show that if \(A \cap B = \{0\}\) , then A - B is a closed, locally compact, cone not containing any line (method as in II, p.~67, exerc. 16). Deduce that there exists a closed hyperplane that supports both A and B, that separates A from B and such that \(H \cap A = H \cap B = \{0\}\) .

\begin{enumerate}
\def\labelenumi{\alph{enumi})}
\setcounter{enumi}{2}
\tightlist
\item
  Give an example of a locally compact closed convex cone A such that \(\mathbf{A} - \mathbf{A}\) is not locally compact (cf.~II, p.~78, exerc. 11).
\end{enumerate}

\(d)\) Let \(\mathbf{A}\) be a locally compact closed convex set in \(\mathbf{E}\) , which does not contain a line. Let \(x_0\) be a point of \(\mathbf{A}\) , \(\mathbf{C}_{\mathbf{A}}\) the asymptotic cone of \(\mathbf{A}\) (cf.~II, p.~67, exerc. 14) and \(\mathbf{H}\) a support hyperplane of \(x_0 + \mathbf{C}_{\mathbf{A}}\) passing through \(x_0\) and such that \((x_0 + \mathbf{C}_{\mathbf{A}}) \cap \mathbf{H} = \{x_0\}\) . If \(f(x) = a\) is the equation of \(\mathbf{H}\) and if \(f(x) \geqslant a\) in \(x_0 + \mathbf{C}_{\mathbf{A}}\) , then show that for every real number \(b\) , the set of the \(y \in \mathbf{A}\) such that \(f(y) \leqslant b\) is compact.

¶ 22) By an extremal ray of a convex subset A of a vector space E we mean a closed half line D contained in A, such that, for all \(x \in \mathbf{D}\) and every open segment with end points \(a, b\) in A, and which contains \(x\) , it is necessarily the case that \(a \in \mathbf{D}\) and \(b \in \mathbf{D}\) ; the end point of D is an extremal point of A.

\begin{enumerate}
\def\labelenumi{\alph{enumi})}
\item
  In a Hausdorff locally convex space E, show that every locally compact closed convex set, not containing a line, is the closed convex envelope of the union of its extremal rays and its extremal points. (Suppose the contrary, and writing B for this closed convex envelope, note first that by exerc. 21, d), there exists a closed hyperplane H so that H ∩ A is compact and non-empty and H ∩ B = ∅. Show then that if \(a \in H \cap A\) is an extremal point of H ∩ A (therefore not an extremal point of A by hypothesis) and if the open segment S with end points b, c contained in A and not contained in H, contains a, then the line D containing S necessarily contains a segment containing a and whose end points are extremal points of A, or contains an extremal ray of A containing a.)
\item
  Prove that if E is finite dimensional, then every closed convex set in E, that does not contain any line, is the convex envelope of the union of its extremal points and its extremal rays (argue by induction on the dimension of E).
\end{enumerate}

\begin{enumerate}
\def\labelenumi{\arabic{enumi})}
\setcounter{enumi}{22}
\tightlist
\item
  In \(R^{3}\) , consider a closed convex set A with an interior point, whose frontier F contains two open segments S, T lying in two non-parallel lines D, \(D'\) (the points of S, T are thus non-extremal in A), and all of whose other frontier points are extremal (one shows how to define such convex sets). For every \(x \in R^{3}\) , put \(f(x) = (d(x, \mathbf{D}))^{2}\) . Let B the convex closed set in \(R^{4} = R^{3} \times R\) formed by the pairs \((x, \zeta)\) such that \(x \in A\) and \(\zeta \geqslant f(x)\) . Show that in B the set of extremal points, the union of the extremal rays, the set of end points of extremal rays, and all the unions of two or three of these sets are not closed and non-empty.
\end{enumerate}

¶ 24) In \(\mathbf{R}^n\) , every intersection of finitely many closed half spaces (resp. of closed half-spaces determined by hyperplanes passing through the same point) is called a polyhedron (resp. a polyhedral cone). A convex set \(\mathbf{C} \subset \mathbf{R}^n\) is locally polyhedral in a point \(x \in \mathbf{C}\) if there is a neighbourhood \(\mathbf{V}\) of \(x\) in \(\mathbf{C}\) which is a polyhedron.

\begin{enumerate}
\def\labelenumi{\alph{enumi})}
\item
  Show that, if a closed convex set \(\mathbf{C} \subset \mathbf{R}^n\) is locally polyhedral at a point \(x \in \mathbf{C}\) , then the cone with vertex \(x\) generated by \(\mathbf{C}\) is polyhedral.
\item
  Show that a compact convex set in \(\mathbf{R}^n\) that is locally polyhedral at each of its points is a polyhedron (use \(a\) ).
\item
  Let \(\mathbf{P} \subset \mathbf{R}^n\) be a closed convex set with an interior point. Show that the following conditions are equivalent:
\end{enumerate}

α) P is a polyhedron.

β) P has only a finite number of facets (II, p.~87, exerc. 3).

\(\gamma)\) P is the convex envelope of a set which is the union of finitely many points and finitely many closed half lines.

(To show that \(\alpha\) ) implies \(\beta\) ) take P as the intersection of the smallest possible number of closed half spaces, and show that the hyperplanes defining these half-spaces are generated by facets of dimension \(n - 1\) of P. To show that \(\beta\) ) implies \(\gamma\) ), argue by induction on \(n\) . Finally, to see that \(\gamma\) ) implies \(\alpha\) ), consider the polar \(\mathrm{P}^{\circ}\) of P.)

\begin{enumerate}
\def\labelenumi{\alph{enumi})}
\setcounter{enumi}{3}
\item
  Show that every convex polyhedron \(\mathbf{P}\) can be written in the form \(\mathrm{Q} + \mathrm{C}_{\mathrm{P}}\) , where \(\mathbf{Q}\) is a compact polyhedron and \(\mathrm{C}_{\mathrm{P}}\) the asymptotic cone of \(\mathbf{P}\) . A non compact polyhedron cannot be parabolic.
\item
  Show that every facet of a convex polyhedron is an ultrafacet (II, p.~88, exerc. 7, d)) (argue by induction on n).
\end{enumerate}

¶ 25) a) Let \(\mathbf{C} \subset \mathbf{R}^n\) be a closed convex cone with vertex 0. Show that the projections of C on every 2-dimensional subspace of \(\mathbf{R}^n\) are closed if and only if C is a polyhedral cone (exerc. 24). (Reduce to the case when C contains no line : argue by induction on \(n\) using the existence of a compact sole S of C (II, p.~90, exerc. 21, a)), and project onto a hyperplane parallel to a line joining 0 to an extremal point of S, and deduce that S is locally polyhedral (exerc. 24). b) Deduce from a) that, if we give \(\mathbf{R}^n\) the order for which C is the set of elements \(\geqslant 0\) , then, every positive linear form on any vector subspace F of \(\mathbf{R}^n\) can be extended to a positive linear form on \(\mathbf{R}^n\) if, and only if, C is a polyhedral cone (apply a) to the polar cone \(\mathbf{C}^\circ\) ). If C is the cone in \(\mathbf{R}^3\) generated by the \((\xi_1, \xi_2, \xi_3)\) such that \(\xi_1 = 1\) , \(\xi_3 \geqslant (\xi_2^3)^-\) , and F is the subspace \(\xi_3 = 0\) give an example of a positive linear form on F that cannot be extended to a positive linear form on \(\mathbf{R}^3\) .

\begin{enumerate}
\def\labelenumi{\alph{enumi})}
\setcounter{enumi}{2}
\tightlist
\item
  Let A be a polyhedron in \(\mathbf{R}^n\) . Then, the convex envelope of \(\mathbf{A} \cup \mathbf{B}\) is closed for every polyhedron B, if and only if, A is compact (use exerc. 24, d)).
\end{enumerate}

\begin{enumerate}
\def\labelenumi{\arabic{enumi})}
\setcounter{enumi}{25}
\tightlist
\item
  \begin{enumerate}
  \def\labelenumii{\alph{enumii})}
  \tightlist
  \item
    Let E be a Hausdorff, locally convex space and let A be a cap of a convex set C in E. If \(s \in C\) is a point not belonging to A and if B is the cone with vertex s generated by A show that the closure of \(\mathbf{B} \cap (\mathbf{C} \cap \complement\mathbf{A})\) is a cap in B.
  \end{enumerate}
\end{enumerate}

\begin{enumerate}
\def\labelenumi{\alph{enumi})}
\setcounter{enumi}{1}
\item
  Suppose that E is finite dimensional. Show that every cap A of a closed convex set C in E can be obtained in the following manner: consider a facet F of C (II, p.~87, exerc. 3) and a hyperplane H in the affine linear variety V generated by F, such that F is entirely on one side of H, take as A the set of points of F contained between H and a hyperplane H' of F parallel to H (use a) and prop. 4 of II, p.~38). Every extremal point of a facet of C is an extremal point of C.
\item
  Give an example of a compact convex set C in a Hausdorff locally convex space E and of a cap A of C such that A and \(C \cap \complement A\) each generate E and that A and \(C \cap \complement A\) cannot be separated by a closed hyperplane of E (cf.~II, p.~78, exerc. 11).
\end{enumerate}

\begin{enumerate}
\def\labelenumi{\arabic{enumi})}
\setcounter{enumi}{26}
\item
  Let C be a closed convex set in a product of lines, E, and let \(a\) be an extremal point of C. Show that for every neighbourhood V of \(a\) in C, there exists an open half-space F in E such that \(a \in \mathrm{F} \cap \mathrm{C} \subset \mathrm{V}\) . (Reduce to the case when C is compact.)
\item
  Let I be a non enumerable infinite set. Show that every cap of the cone \(R_{+}^{I}\) in \(R^{I}\) is contained in the sum of subspaces of the form \(R^{J}\) , where J is an enumerable subset of I (use prop. 4 of II, p.~38). Deduce that there are points of \(R_{+}^{I}\) which do not belong to any cap of \(R_{+}^{I}\) , even though \(R_{+}^{I}\) is the convex closed envelope of the union of its extremal generators.
\item
  \begin{enumerate}
  \def\labelenumii{\alph{enumii})}
  \tightlist
  \item
    Let \((\mathrm{E}_n)\) be a sequence of Hausdorff locally convex spaces and \(\mathrm{E} = \prod_{n}\mathrm{E}_{n}\) their product.
  \end{enumerate}
\end{enumerate}

In each \(E_{n}\) , let \(C_{n}\) be a convex cone with vertex 0 and \(A_{n}\) a cap of \(C_{n}\) . Show that there exists a cap of \(C = \prod C_{n}\) which contains \(\prod A_{n}\) (argue as in prop. 5 of II, p.~59).

\begin{enumerate}
\def\labelenumi{\alph{enumi})}
\setcounter{enumi}{1}
\tightlist
\item
  Let \((\overline{\mathbf{E}}_n, \phi_{nm})\) be an enumerable directed projective system of Hausdorff locally convex spaces and let \(\mathrm{E} = \varprojlim \mathrm{E}_n\) be its projective limit. For all \(n\) , let \(\mathbf{C}_n\) be a convex cone of vertex 0 such that \((\mathbf{C}_n)\) is a projective system of sets. Show that if, for each \(n\) , the set \(C_n\) is the union of its caps, then this is also true of \(\mathrm{C} = \varprojlim \mathrm{C}_n\) (use \(a\) ). In particular, if the \(\mathbf{C}_n\) are cones with compact soles then \(\mathbf{C}\) in the closed convex envelope of the union of its extremal generators.
\end{enumerate}

\begin{enumerate}
\def\labelenumi{\arabic{enumi})}
\setcounter{enumi}{29}
\tightlist
\item
  Let \(\mathbf{E}\) be a Hausdorff locally convex space and \(\mathbf{A}\) a cap of \(\mathbf{C}\) , a closed convex set in \(\mathbf{E}\) . Show that if \(a \in \mathbf{A}\) is an extremal point of \(\mathbf{A}\) then the facet \(\mathbf{F}\) of \(a\) in \(\mathbf{C}\) (II, p.~87, exerc. 3) is of dimension \(\leqslant 1\) (use exerc. 26 of II, p.~91). Deduce that \(\mathbf{F} \cap \mathbf{A}\) is a cap of \(\mathbf{C}\) .
\end{enumerate}

¶ * 31) Let X be the compact interval \([0, 1]\) of R and \(\Phi\) the set formed from the continuous real valued functions defined in X and the functions \(t \mapsto |t - a|^{-\alpha}\) , where \(a \in X\) and \(0 < \alpha < 1\) (we put \(0^{-\alpha} = +\infty\) for \(\alpha > 0\) ). In the space \(\mathcal{M}(X)\) of measures on X, let \(M_{\Phi}^{+}\) be the set of the measures \(\mu \geqslant 0\) such that all the functions of \(\Phi\) are \(\mu\) -integrable.

\begin{enumerate}
\def\labelenumi{\alph{enumi})}
\item
  Give \(\mathcal{M}_{\Phi}^{+}\) the uniform structure induced by the product structure of \(\mathbf{R}^{\Phi}\) . Show that \(\mathcal{M}_{\Phi}^{+}\) is a proper convex complete cone for this uniform structure. (Note that for every function \(f \in \Phi\) there exists \(g \in \Phi\) such that, for all \(\varepsilon > 0\) , there exists \(u \in \mathcal{C}(\mathbf{X}; \mathbf{R})\) such that \(0 \leqslant f - u \leqslant \varepsilon g\) . b) Show that the cone \(\mathcal{M}_{\Phi}^{+}\) has no extremal generator. (Observe that if \(\mu \in \mathcal{M}_{\Phi}^{+}\) , then all the measures \(\lambda\) such that \(0 \leqslant \lambda \leqslant \mu\) belong to \(\mathcal{M}_{\Phi}^{+}\) .)
\item
  Show that the set S of the \(\mu \in \mathcal{M}_{\Phi}^{+}\) such that \(\mu(1) = 1\) is a sole of the cone \(\mathcal{M}_{\Phi}^{+}\) and a simplex in \(\mathbf{R}^{\Phi}\) (II, p.~71, exerc. 41). *
\end{enumerate}

\begin{enumerate}
\def\labelenumi{\arabic{enumi})}
\setcounter{enumi}{31}
\tightlist
\item
  Let \(\mathbf{E}\) and \(\mathbf{F}\) be two Hausdorff locally convex spaces, let \(\mathbf{A}\) be a convex subset of \(\mathbf{E}\) , and \(u\) a linear mapping of \(\mathbf{E}\) in \(\mathbf{F}\) .
\end{enumerate}

\begin{enumerate}
\def\labelenumi{\alph{enumi})}
\item
  The inverse image under \(u\) of a support variety of \(u(\mathbf{A})\) (II, p.~87, exerc. 3, c)) is a support variety of \(\mathbf{A}\) .
\item
  If \(\mathbf{A}\) is compact and \(u\) is continuous, then every extremal point of \(u(\mathbf{A})\) is the image under \(u\) of an extremal point of \(\mathbf{A}\) .
\item
  If \(\mathbf{A}\) is a locally compact cone with vertex 0 and if \(u\) is continuous then every extremal generator of \(u(\mathbf{A})\) is the image under \(u\) of an extremal generator of \(\mathbf{A}\) .
\end{enumerate}

¶ 33) Let E be a Hausdorff locally convex space and A a subset of E.

\begin{enumerate}
\def\labelenumi{\alph{enumi})}
\item
  Denote by \(\Gamma_0(\mathbf{A})\) the set of points \(x \in \mathbf{E}\) such that, for every continuous linear mapping \(u\) of \(\mathbf{E}\) in a finite dimensional vector space, the image \(u(x)\) belongs to the convex envelope of \(u(\mathbf{A})\) . This comes to the same as saying that for every closed linear variety \(\mathbf{V}\) of \(\mathbf{E}\) containing \(x\) and of finite codimension \(n > 0\) , there exists a subset of \(\mathbf{A}\) having at most \(n + 1\) elements, and of which the convex envelope meets \(\mathbf{V}\) . Show that \(\Gamma_0(\mathbf{A})\) is a convex set containing \(\mathbf{A}\) , that \(\Gamma_0(\Gamma_0(\mathbf{A})) = \Gamma_0(\mathbf{A})\) and that \(\Gamma_0(\mathbf{A})\) is contained in the closed convex envelope of \(\mathbf{A}\) (use prop. 4 of II, p.~38).
\item
  Let \((x_{\alpha})_{\alpha \in \mathbb{I}}\) be a family of elements of A and \((\lambda_{\alpha})_{\alpha \in \mathbb{I}}\) a family of positive numbers such that \(\sum_{\alpha \in \mathbb{I}}\lambda_{\alpha} = 1\) and that the family \((\lambda_{\alpha}x_{\alpha})\) is summable in E. Show that the sum \(s = \sum_{\alpha \in \mathbb{I}}\lambda_{\alpha}x_{\alpha}\) belongs
\end{enumerate}

to \(\Gamma_0(\mathbf{A})\) . (With the aid of \(a\) ) reduce to the case where \(\mathbf{E}\) is of finite dimension and identical with the linear variety generated by the \(\lambda_{\alpha}x_{\alpha}\) ; then argue by reductio ad absurdum, considering, for every finite subset \(\mathbf{J}\) of \(\mathbf{I}\) , a closed hyperplane \(\mathbf{H}_{\mathbf{J}}\) that passes through \(s\) and does not meet the convex envelope of the set of the \(x_{\alpha}\) such that \(\alpha \in \mathbf{J}\) , then using the compactness of the unit sphere in a finite dimensional space.)

\begin{enumerate}
\def\labelenumi{\alph{enumi})}
\setcounter{enumi}{2}
\item
  Show that if A is compact, then \(\Gamma_0(A)\) is identical with the convex closed envelope of A. d) If K is a compact convex set in E, and A the set of its extremal points show that \(K = \Gamma_0(A)\) (use exerc. 22, b) of II, p.~90, and exerc. 32).
\item
  With the notations of II, p.~74, exerc. 3, let A be the set formed of the \(\varepsilon_{x}\) , where \(x\) varies in the set of rational numbers such that \(0 \leqslant x \leqslant 1\) . Show that \(\Gamma_0(A)\) is distinct from the convex envelope of A and from the convex closed envelope of A.
\end{enumerate}

¶ 34) Let S be a closed convex set of E, a Hausdorff locally convex space, and A a subset of S such that \(\mathbf{S} = \Gamma_0(\mathbf{A})\) (exerc. 33), and let \(\mathbf{S}_0\) be the convex envelope of A (so that \(\mathbf{S} = \bar{\mathbf{S}}_0\) ). Let N be a closed convex subset of E containing a closed linear variety of finite codimension, \(\mathbf{M} = \mathbf{S} \cap \mathbf{N}\) and \(\mathbf{M}_0 = \mathbf{S}_0 \cap \mathbf{N}\) .

\begin{enumerate}
\def\labelenumi{\alph{enumi})}
\item
  Show that \(\mathbf{M} = \overline{\mathbf{M}}_0\) . (Note, using exerc. 33, a), that every closed linear variety of finite codimension in E, containing a point \(x \in \mathbf{M}\) , meet \(\mathbf{M}_0\) , and use the prop. 4 of II, p.~38.)
\item
  Suppose that for every finite subset F of A, the intersection of N and of the facet (in S) of each point of the convex envelope of F, is compact or of finite dimension and does not contain a line. Show then that M is the convex closed envelope of the set of its extremal points. (By the aid of a), this reduces to proving that every point of \(M_{0}\) is contained in the convex closed envelope of the set of extremal points of M. Use exerc. 3, e) of II, p.~87, the Krein-Milman th. (II, p.~55) and exerc. 22, b) of II, p.~90.) Deduce that every closed support hyperplane of M contains an extremal point of M.
\end{enumerate}

\begin{enumerate}
\def\labelenumi{\arabic{enumi})}
\setcounter{enumi}{34}
\tightlist
\item
  Let I be a non enumerable set, write \(E = R^{(I)}\) and \(E' = E \times R\) . Denote the canonical basis of E by \((e_{\alpha})_{\alpha \in I}\) and let s be the element \((0, 1)\) of \(E'\) . Define a separating duality between E and \(E'\) , by \(\langle e_{\alpha}, e_{\beta} \rangle = \delta_{\alpha\beta}, \langle e_{\alpha}, s \rangle = 1\) for all \(\alpha \in I\) . Let C be the pointed cone \(R_{+}^{I}\) in E. a) Show that the topologies induced on C by \(\sigma(E, E')\) and by the norm \(p(x) = \sum_{\alpha \in I} |x_{\alpha}|\) on E coincide.
\end{enumerate}

\begin{enumerate}
\def\labelenumi{\alph{enumi})}
\setcounter{enumi}{1}
\tightlist
\item
  Show that the uniform structure induced on C by \(\sigma(\mathrm{E},\mathrm{E}^{\prime})\) is not metrisable.
\end{enumerate}

\begin{enumerate}
\def\labelenumi{\arabic{enumi})}
\setcounter{enumi}{35}
\tightlist
\item
  Consider the space \(\mathrm{E} = \mathbf{R}^{(\mathbf{N})}\) , with the weak topology \(\sigma (\mathbf{R}^{(\mathbf{N})},\mathbf{R}^{\mathbf{N}})\) ; let C be the closed convex cone in E formed by the points \(x = (x_{n})\) such that \(x_{n}\geqslant 0\) for all \(n\) .\\
\end{enumerate}

\begin{enumerate}
\def\labelenumi{\alph{enumi})}
\item
  Let \(x = (x_{n})\) be a point of C and let J be the finite set of integers \(n\) for which \(x_{n} > 0\) ; if \(m\) is the number of elements of J, then let A be the set of the points \(y = (y_{n})\) of C such that \(y_{n} = 0\) for \(n\in \mathrm{J}\) and \(\sum_{k\in J}y_kx_k^{-1}\leqslant m\) . Show that A is a cap of C containing \(x\) .
\item
  Show that there does not exist a cap B in C such that C is the union of the sets nB for n \textgreater{} 0. (Let p be the restriction of the gauge of B to C; p will be finite in C and, if \((e_{n})\) is the canonical basis of E, we have \(p(e_{n}) > 0\) for all n (II, p.~58, prop. 4), and the points \(z^{(n)} = e_{n}/p(e_{n})\) belong to B; but show that there exists \(z' \in R^{N}\) such that the sequence \((\langle z^{(n)}, z' \rangle)\) is not bounded.)
\end{enumerate}

\begin{enumerate}
\def\labelenumi{\arabic{enumi})}
\setcounter{enumi}{36}
\tightlist
\item
  Let F be the Banach space \(l^{1}(\mathbf{N})\) of summable sequences \(x = (x_{n})\) of real numbers and let E be the space of sequences \(y = (y_{n})\) which tend to 0; we give F the weak topology \(\sigma(\mathrm{F}, \mathrm{E})\) where E and F are in separating duality using the form \(\mathbf{B}(x, y) = \sum x_{n} y_{n}\) .
\end{enumerate}

\begin{enumerate}
\def\labelenumi{\alph{enumi})}
\item
  Let \(\mathbf{C}\) be the convex cone in \(\mathbf{F}\) formed of the points \(x = (x_{n})\) such that \(x_{n} \geqslant 0\) for all \(n\) . Show that \(\mathbf{C}\) is closed in \(\mathbf{F}\) .
\item
  Let A be the set of the \(x = (x_{n}) \in \mathbf{C}\) such that \(\sum_{n} x_{n} \leqslant 1\) . Show that A is a cap of C, that is metrisable for the topology induced by that of F, and that C is the union of the sets nA for \(n > 0\) .
\item
  Show that a sole of C is not compact (such a set S would be the set of the \(x = (x_{n}) \in \mathbf{C}\) such that \(\sum_{n} z_{n} x_{n} = 1\) , where \((z_{n}) \in \mathbf{E}\) and \(z_{n} > 0\) for all n.~If \(e_{m} = (\delta_{mn})_{n \geqslant 0}\) , the points \(z_{n}^{-1} e_{n}\) belong to S, but do not form a relatively compact set in F).
\end{enumerate}

\(d\) ) Show that \(\mathbf{C}\) is not metrisable for the topology induced by that of \(\mathbf{F}\) (use Baire's th., noting that there is no point in A that is an interior point relative to the subspace C).

\begin{enumerate}
\def\labelenumi{\arabic{enumi})}
\setcounter{enumi}{37}
\tightlist
\item
  Let E be a Hausdorff locally convex space and X a compact convex set in E. Denote by \(\mathcal{A}(X)\) the set of continuous affine functions in X (not necessarily restrictions to X of continuous affine functions in E, cf.~II, p.~78, exerc. 11). For every real valued function f that is bounded above in X, put \(\tilde{f}(x) = \inf_{h}(h(x))\) for all \(x \in X\) where h varies in the set of function of \(\mathcal{A}(X)\) such that \(h \geqslant f\) .
\end{enumerate}

\begin{enumerate}
\def\labelenumi{\alph{enumi})}
\item
  Show that \(\tilde{f}\) is an upper semi-continuous concave function. If \(f\) itself is upper semicontinuous and concave, then \(\tilde{f} = f\) (cf.~II, p.~39, prop. 5).
\item
  Suppose that \(f\) is upper semi-continuous. Show that \(\tilde{f}\) is the lower envelope of the function \(\tilde{g}\) , when \(g\) varies in a set of functions that are continuous in \(X\) of which \(f\) is the lower envelope.
\item
  For all \(x \in \mathbf{X}\) , write \(\mathcal{M}_x\) for the set of finite families \(\mu = ((\mu_j, x_j))\) where the \(x_j\) are points of \(\mathbf{X}\) , the \(\mu_j\) are numbers \(\geqslant 0\) satisfying \(\sum_j \mu_j = 1\) , such that \(x = \sum_j \mu_j x_j\) . For each real valued function \(f\) bounded above in \(\mathbf{X}\) , put \(f'(x) = \sup_{\mu \in \mathcal{M}_x} \sum \mu_j f(x_j)\) for all \(x \in \mathbf{X}\) . Show that \(f'\) is a concave function in \(\mathbf{X}\) and that \(f' \leqslant \tilde{f}\) .
\item
  Suppose that \(f\) is continuous in \(\mathbf{X}\) . Given \(\varepsilon > 0\) , let \((\mathrm{U}_k)_{1 \leqslant k \leqslant N}\) be a covering of \(\mathbf{X}\) by convex open sets such that the relations \(x \in \mathrm{U}_k\) , \(y \in \mathrm{U}_k\) imply the inequality \(|f(x) - f(y)| \leqslant \varepsilon\) . Put \(\mathrm{A}_1 = \mathrm{U}_1\) and, for \(k > 1\) , \(\mathrm{A}_k = \mathrm{U}_k \cap \complement(\mathrm{U}_1 \cup \mathrm{U}_2 \cup \mathrm{U}_3 \cup \ldots \cup \mathrm{U}_{k-1})\) . Show that, for all \(x \in \mathbf{X}\) , there exists a family \(\mu = ((\mu_k, x_k))\) of \(N\) terms belonging to \(\mathcal{M}_x\) with \(x_k \in \mathrm{U}_k\) for \(1 \leqslant k \leqslant N\) such that \(\sum_k \mu_k f(x_k) \geqslant f'(x) - 2\varepsilon\) (if we have the inequality \(\sum_j \lambda_j f(y_j) \geqslant f'(x) - \varepsilon\) , group the \(y_j\) belonging to the same \(A_k\) together).
\item
  Deduce from \(d\) that when \(f\) is continuous then \(f'\) is upper semi-continuous and \(f' = \tilde{f}\) . (If \(\mathfrak{U}\) is an ultrafilter on \(\mathbf{X}\) finer than the filter of the neighbourhoods of a point \(x \in \mathbf{X}\) and if \(f'(y) \geqslant r\) for all the points \(y\) of a set belonging to \(\mathfrak{U}\) , show that \(f'(x) \geqslant r - 2\varepsilon\) , by making a family \(\mu_y \in \mathcal{M}_y\) correspond to each \(y\) , which satisfies condition \(d\) ) and proceeding to the limit through \(\mathfrak{U}\) .)
\end{enumerate}

\begin{enumerate}
\def\labelenumi{\arabic{enumi})}
\setcounter{enumi}{38}
\tightlist
\item
  Let \(\mathbf{H}\) be a closed hyperplane in a Hausdorff locally convex space \(\mathbf{E}\) , that does not contain 0 and let \(S\) be a compact simplex contained in \(\mathbf{H}\) (II, p.~71, exerc. 41).
\end{enumerate}

\begin{enumerate}
\def\labelenumi{\alph{enumi})}
\item
  Let C be the cone with vertex 0 generated by S. Show that if \((x_{i})_{i\in \mathbf{I}}\) and \((y_{j})_{j\in \mathbf{J}}\) are two finite families of points of C such that \(\sum_{i\in \mathbf{I}}x_i = \sum_{j\in \mathbf{J}}y_j\) , then there exists a finite family \((z_{ij})_{(i,j)\in \mathbf{I}\times \mathbf{J}}\) of points of C such that \(x_{i} = \sum_{j\in \mathbf{J}}z_{ij}\) for all \(i\in \mathbf{I}\) and \(y_{j} = \sum_{i\in \mathbf{I}}z_{ij}\) for all \(j\in \mathbf{J}\) (argue by induction to reduce to the case \(\mathrm{I} = \mathrm{J} = \{1,2\}\) ).
\item
  Let \(f\) be a convex, upper semi-continuous function in S. Show that the function \(\tilde{f}\) (defined in exerc. 38) is an affine function. (First reduce to the case where \(f\) is continuous by using exerc. 38, b) and II, p.~81, exerc. 28. Next use the fact if \(f\) is continuous, then \(\tilde{f} = f'\) (exerc. 38, e)), and show that \(f'\) is convex using a) to bound \(f'(\alpha_1 x_1 + \alpha_2 x_2)\) above when \(\alpha_1 \geqslant 0\) , \(\alpha_2 \geqslant 0\) , \(\alpha_1 + \alpha_2 = 1\) .)
\end{enumerate}

\begin{enumerate}
\def\labelenumi{\arabic{enumi})}
\setcounter{enumi}{39}
\tightlist
\item
  Let \(\mathbf{X}\) be a compact convex set in \(\mathbf{E}\) , a Hausdorff locally convex space, and let \(f\) be an upper semi-continuous function that is bounded below. Let \(g\) be a lower semi-continuous concave function, such that \(g \geqslant f\) . Show (with that notations of exerc. 38) that \(g \geqslant \tilde{f}\) . (Reduce to the case where \(\inf_{x \in \mathbf{X}} (g(x) - f(x)) > 0\) . If \((f_{\alpha})\) is a decreasing directed family of continuous
\end{enumerate}

functions such that \(f = \inf(f_{\alpha})\) , show that then also \(\inf_{x \in X} (g(x) - f_{\alpha}(x)) \geqslant 0\) for \(\alpha \geqslant \alpha_0\) , and so reduce the problem to the case where \(f\) is continuous. Then use exerc. 38, e).

\begin{enumerate}
\def\labelenumi{\arabic{enumi})}
\setcounter{enumi}{40}
\tightlist
\item
  Let S be a compact simplex contained in a Hausdorff locally convex space E (II, p.~71 exerc. 41), and \(f\) an upper semi-continuous convex function that is bounded below. Let \(g\) be a lower semi-continuous function that is concave and such that \(g \geqslant f\) .
\end{enumerate}

\begin{enumerate}
\def\labelenumi{\alph{enumi})}
\item
  If we write \(u = \tilde{f}\) , \(v = -(-g)\tilde{f}\) , then \(u\) and \(v\) are affine functions such that \(u \leqslant v\) (use exerc. 39, b) and exerc. 40).
\item
  Show that there exists an affine function \(h\) , continuous in \(\mathbf{X}\) and such that \(f \leqslant h \leqslant g\) (D. Edwards' th.). (We can replace \(f\) by \(u\) and \(g\) by \(v\) . Construct three sequences \((u_m), (v_m), (h_m)\) of affine functions such that in \(\mathbf{X}\) , the function \(u_m\) is upper semi-continuous, the function \(v_m\) is lower semi-continuous and the function \(h_m\) is continuous and
\end{enumerate}

\[
u - \frac {1}{2 ^ {m}} \leqslant u _ {m} <   h _ {m} <   v _ {m} \leqslant v + \frac {1}{2 ^ {m}} \text { and } \| h _ {n + 1} - h _ {n} \| \leqslant \frac {1}{2 ^ {n + 1}}.
\]

Use exerc. 29 of II, p.~81, for this.)

